\documentclass[12pt,a4paper]{ctexbook}
\usepackage{geometry}
\geometry{margin=2.2cm}
\usepackage{setspace}
\setstretch{1.35}
\usepackage{hyperref}
\title{花间集·huajianji-3-juan}
\author{古典文献整理}
\date{}
\begin{document}
\maketitle
\tableofcontents
\newpage
\section{谒金门·春漏促}
\textbf{韦庄}\par\vspace{0.5em}
春漏促，金烬暗挑残烛。\par
一夜帘前风撼竹，梦魂相断续。\par
有个娇娆如玉，夜夜绣屏孤宿。\par
宋袍琵琶寻旧曲，远山眉黛绿。\par

\section{谒金门·空相忆}
\textbf{韦庄}\par\vspace{0.5em}
空相忆，无计得传消息。\par
天上嫦娥人不识，寄书何处觅。\par
新睡觉来无力，不忍把伊书迹。\par
满院落花春寂寂，断肠芳草碧。\par

\section{江城子·恩重娇多情易伤}
\textbf{韦庄}\par\vspace{0.5em}
恩重娇多情易伤，漏更长，解鸳鸯。\par
朱唇未动，先觉口脂香。\par
缓揭绣衾抽皓腕，移凤枕，枕潘郎。\par

\section{江城子·髻鬟狼藉黛眉长}
\textbf{韦庄}\par\vspace{0.5em}
髻鬟狼藉黛眉长，出兰房，别檀郎。\par
角声呜咽，星斗渐微茫。\par
露冷月残人未起，留不住，泪千行。\par

\section{河传·何处}
\textbf{韦庄}\par\vspace{0.5em}
何处，烟雨，隋堤春暮。\par
柳色葱笼，画桡金缕，翠旗高颭香风，水光融。\par
青娥殿脚春妆媚，轻云里，绰约司花妓。\par
江都宫阙，清淮月映迷楼，古今愁。\par

\section{河传·春晚}
\textbf{韦庄}\par\vspace{0.5em}
春晚，风暖，锦城花满。\par
狂杀游人，玉鞭金勒，寻胜驰骤轻尘，惜良晨。\par
翠娥争劝临邛酒，纤纤手，拂面垂丝柳。\par
归时烟里，钟鼓正是黄昏，暗销魂。\par

\section{河传·锦浦}
\textbf{韦庄}\par\vspace{0.5em}
锦浦，春女，绣衣金缕。\par
雾薄云轻，花深柳暗，时节正是清明，雨初晴。\par
玉鞭魂断烟霞路，莺莺语，一望巫山雨。\par
香尘隐映，遥见翠槛红楼，黛眉愁。\par

\section{天仙子·怅望前回梦里期}
\textbf{韦庄}\par\vspace{0.5em}
怅望前回梦里期，看花不语苦寻思。\par
露桃花里小腰肢，眉眼细，鬓云垂，唯有多情宋玉知。\par

\section{天仙子·深夜归来长酩酊}
\textbf{韦庄}\par\vspace{0.5em}
深夜归来长酩酊，扶入流苏犹未醒。\par
醺醺酒气麝兰和，惊睡觉，笑呵呵，长道人生能几何。\par

\section{天仙子·蟾彩霜华夜不分}
\textbf{韦庄}\par\vspace{0.5em}
蟾彩霜华夜不分，天外鸿声枕上闻，绣衾香冷懒重熏。\par
人寂寂，叶纷纷，才睡依前梦见君。\par

\section{天仙子·梦觉云屏依旧空}
\textbf{韦庄}\par\vspace{0.5em}
梦觉云屏依旧空，杜鹃声咽隔帘栊，玉郎薄倖去无踪。\par
一日日，恨重重，泪界莲腮两线红。\par

\section{天仙子·金似衣裳玉似身}
\textbf{韦庄}\par\vspace{0.5em}
金似衣裳玉似身，眼如秋水鬓如云，霞裙月帔一群群。\par
来洞口，望烟分，刘阮不归春日曛。\par

\section{喜迁莺·人汹汹}
\textbf{韦庄}\par\vspace{0.5em}
人汹汹，鼓冬冬，襟袖五更风。\par
大罗天上月朦胧，骑马上虚空。\par
香满衣，云满路，鸾凤绕身飞舞。\par
霓旌绛节一群群，引见玉华君。\par

\section{喜迁莺·街鼓动}
\textbf{韦庄}\par\vspace{0.5em}
街鼓动，禁城开，天上探人回，凤衔金榜出云来，平地一声雷。\par
莺已迁，龙已化，一夜满城车马。家家楼上簇神仙，争看鹤冲天。\par

\section{思帝乡·云髻坠}
\textbf{韦庄}\par\vspace{0.5em}
云髻坠，凤钗垂。\par
髻坠钗垂无力，枕函欹。\par
翡翠屏深月落，漏依依。\par
说尽人间天上，两心知。\par

\section{思帝乡·春日游}
\textbf{韦庄}\par\vspace{0.5em}
春日游，杏花吹满头。\par
陌上谁家年少，足风流。\par
妾拟将身嫁与，一生休。\par
纵被无情弃，不能羞！\par

\section{诉衷情·烛尽香残帘半卷}
\textbf{韦庄}\par\vspace{0.5em}
烛尽香残帘半卷，梦初惊。\par
花欲谢，深夜，月胧明。\par
何处按歌声，轻轻。\par
舞衣尘暗生，负春情。\par

\section{诉衷情·碧沼红芳烟雨静}
\textbf{韦庄}\par\vspace{0.5em}
碧沼红芳烟雨静，倚兰桡。\par
垂玉珮，交带，袅纤腰。\par
鸳梦隔星桥，迢迢。\par
越罗香暗销，坠花翘。\par

\section{上行杯·芳草灞陵春岸}
\textbf{韦庄}\par\vspace{0.5em}
芳草灞陵春岸，柳烟深，满楼弦管，一曲离声肠寸断。\par
今日送君千万，红楼玉盘金镂盏，须劝珍重意，莫辞满。\par

\section{上行杯·白马玉鞭金辔}
\textbf{韦庄}\par\vspace{0.5em}
白马玉鞭金辔，少年郎，离别容易，迢递去程千万里。\par
惆怅异乡云水，满酌一杯劝和泪，须愧珍重意，莫辞醉。\par

\section{女冠子·四月十七}
\textbf{韦庄}\par\vspace{0.5em}
四月十七，正是去年今日。\par
别君时，忍泪佯低面，含羞半敛眉。\par
不知魂已断，空有梦相随。\par
除却天边月，没人知。\par

\section{女冠子·昨夜夜半}
\textbf{韦庄}\par\vspace{0.5em}
昨夜夜半，枕上分明梦见：语多时，依旧桃花面，频低柳叶眉。\par
半羞还半喜，欲去又依依。觉来知是梦，不胜悲。\par

\section{更漏子·钟鼓寒}
\textbf{韦庄}\par\vspace{0.5em}
钟鼓寒，楼阁暝，月照古桐金井。\par
深院闭，小庭空，落花香露红。\par
烟柳重，春雾薄，灯背水窗高阁。\par
闲倚户，暗沾衣，待郎郎不归。\par

\section{酒泉子·月落星沉}
\textbf{韦庄}\par\vspace{0.5em}
月落星沉，楼上美人春睡。\par
绿云倾，金枕腻，画屏深。\par
子规啼破相思梦，曙色东方才动。\par
柳烟轻，花露重，思难任。\par

\section{木兰花·独上小楼春欲暮}
\textbf{韦庄}\par\vspace{0.5em}
独上小楼春欲暮，愁望玉关芳草路。\par
消息断，不逢人，欲敛细眉归绣户。\par
坐看落花空叹息，罗袂湿斑红泪滴。\par
千山万水不曾行，魂梦欲教何处觅。\par

\section{小重山·一闭昭陽春又春}
\textbf{韦庄}\par\vspace{0.5em}
一闭昭陽春又春，夜寒宫漏永，梦君恩。\par
昏思陈事暗消魂，罗衣湿，红袂有啼痕。\par
歌吹隔重阍，绕庭芳草绿，倚长门。\par
万般惆怅向谁论？凝情立，宫殿欲黄昏。\par

\section{浣溪沙·红蓼渡头秋正雨}
\textbf{薛昭蕴}\par\vspace{0.5em}
红蓼渡头秋正雨，印沙鸥迹自成行，整鬟飘袖野风香。\par
不语含颦深浦里，几回愁煞棹船郎，燕归帆尽水茫茫。\par

\section{浣溪沙·钿匣菱花锦带垂}
\textbf{薛昭蕴}\par\vspace{0.5em}
钿匣菱花锦带垂，静临栏槛卸头时，约鬟低珥算归期。\par
茂苑草青湘渚阔，梦余空有漏依依，二年终日损芳菲。\par

\section{浣溪沙·粉上依稀有泪痕}
\textbf{薛昭蕴}\par\vspace{0.5em}
粉上依稀有泪痕，郡庭花落欲黄昏，远情深恨与谁论。\par
记得去年寒食日，延秋门外卓金轮，日斜人散暗消魂。\par

\section{浣溪沙·握手河桥柳似金}
\textbf{薛昭蕴}\par\vspace{0.5em}
握手河桥柳似金，蜂须轻惹百花心，蕙风兰思寄清琴。\par
意满便同春水满，情深还似酒杯深，楚烟湘月两沉沉。\par

\section{浣溪沙·帘下三间出寺墙}
\textbf{薛昭蕴}\par\vspace{0.5em}
帘下三间出寺墙，满街垂杨绿陰长，嫩红轻翠间浓妆。\par
瞥地见时犹可可，却来闲处暗思量，如今情事隔仙乡。\par

\section{浣溪沙·江馆清秋缆客船}
\textbf{薛昭蕴}\par\vspace{0.5em}
江馆清秋缆客船，故人相送夜开筵，麝烟兰焰簇花钿。\par
正是断魂迷楚雨，不堪离恨咽湘弦，月高霜白水连天。\par

\section{浣溪沙·倾国倾城恨有余}
\textbf{薛昭蕴}\par\vspace{0.5em}
倾国倾城恨有余，几多红泪泣姑苏，倚风凝睇雪肌肤。\par
吴主山河空落日，越王宫殿半平芜，藕花菱蔓满重湖。\par

\section{浣溪沙·越女淘金春水上}
\textbf{薛昭蕴}\par\vspace{0.5em}
越女淘金春水上，步摇云鬓珮鸣璫，渚风江草又清香。\par
不为远山凝翠黛，只应含恨向斜陽，碧桃花谢忆刘郎。\par

\section{喜迁莺·残蟾落}
\textbf{薛昭蕴}\par\vspace{0.5em}
残蟾落，晓钟鸣，羽化觉身轻。\par
乍无春睡有余酲，杏苑雪初晴。\par
紫陌长，襟袖冷，不是人间风景。\par
回看尘土似前生，休羡谷中莺。\par

\section{喜迁莺·金门晓}
\textbf{薛昭蕴}\par\vspace{0.5em}
金门晓，玉京春，骏马骤轻尘。\par
桦烟深处白衫新，认得化龙身。\par
九陌喧，千户启，满袖桂香风细。\par
杏园欢宴曲江滨，自此占芳辰。\par

\section{喜迁莺·清明节}
\textbf{薛昭蕴}\par\vspace{0.5em}
清明节，雨晴天，得意正当年。\par
马骄泥软绵连乾，香袖半笼鞭。\par
花色融，人竞赏，尽是绣鞍朱鞅。\par
日斜无计更留连，归路草和烟。\par

\section{小重山·春到长门春草青}
\textbf{薛昭蕴}\par\vspace{0.5em}
春到长门春草青，玉阶华露滴，月胧明。\par
东风吹断紫箫声，宫漏促，帘外晓啼莺。\par
愁极梦难成，红妆流宿泪，不胜情。\par
手挼裙带绕阶行，思君切，罗幌暗尘生。\par

\section{小重山·秋到长门秋草黄}
\textbf{薛昭蕴}\par\vspace{0.5em}
秋到长门秋草黄，画梁双燕去，出宫墙。\par
玉箫无复理霓裳，金蝉坠，鸾镜掩休妆。\par
忆昔在昭陽，舞衣红绶带，绣鸳鸯。\par
至今犹惹御炉香，魂梦断，愁听漏更长。\par

\section{离别难·宝马晓鞲雕鞍}
\textbf{薛昭蕴}\par\vspace{0.5em}
宝马晓鞲雕鞍，罗帷乍别情难。\par
那堪春景媚，送君千万里。\par
半妆珠翠落，露华寒。\par
红蜡烛，青丝曲，偏能钩引泪阑干。\par
良夜促，香尘绿，魂欲迷，檀眉半敛愁低。\par
未别心先咽，欲语情难说。\par
出芳草，路东西。\par
摇袖立，春风急，樱花杨柳雨凄凄。\par

\section{相见欢·罗襦绣袂香红}
\textbf{薛昭蕴}\par\vspace{0.5em}
罗襦绣袂香红，画堂中。\par
细草平沙番马，小屏风。\par
卷罗幕，凭妆阁，思无穷。\par
暮雨轻烟魂断，隔帘栊。\par

\section{醉公子·慢绾青丝发}
\textbf{薛昭蕴}\par\vspace{0.5em}
慢绾青丝发，光砑吴绫袜。\par
床上小熏笼，韶州新退红。\par
叵耐无端处，捻得从头污。\par
恼得眼慵开，问人闲事来。\par

\section{女冠子·求仙去也}
\textbf{薛昭蕴}\par\vspace{0.5em}
求仙去也，翠钿金篦尽舍，入崖峦。\par
雾卷黄罗帔，云雕白玉冠。\par
野烟溪洞冷，林月石桥寒。\par
静夜松风下，礼天坛。\par

\section{女冠子·云罗雾縠}
\textbf{薛昭蕴}\par\vspace{0.5em}
云罗雾縠，新授明威法籙，降真函。\par
髻绾青丝发，冠抽碧玉簪。\par
往来云过五，去往岛经三。\par
正遇刘郎使，启瑶缄。\par

\section{谒金门·春满院}
\textbf{薛昭蕴}\par\vspace{0.5em}
春满院，叠损罗衣金线。\par
睡觉水晶帘未卷，檐前双语燕。\par
斜掩金铺一扇，满地落花千片。\par
早是相思肠欲断，忍教频梦见。\par

\section{柳枝·解冻风来末上青}
\textbf{牛峤}\par\vspace{0.5em}
解冻风来末上青，解垂罗袖拜卿。\par
无端袅娜临官路，舞送行人过一生。\par

\section{柳枝·吴王宫里色偏深}
\textbf{牛峤}\par\vspace{0.5em}
吴王宫里色偏深，一簇纤条万缕金。\par
不愤钱塘苏小小，引郎松下结同心。\par

\section{柳枝·桥北桥南千万条}
\textbf{牛峤}\par\vspace{0.5em}
桥北桥南千万条，恨伊张绪不相饶。\par
金羁白马临风望，认得杨家静婉腰。\par

\section{柳枝·狂雪随风扑马飞}
\textbf{牛峤}\par\vspace{0.5em}
狂雪随风扑马飞，惹烟无力被春欺。\par
莫教移入灵和殿，宫女三千又妒伊。\par

\section{柳枝·袅翠笼烟拂暖波}
\textbf{牛峤}\par\vspace{0.5em}
袅翠笼烟拂暖波，舞裙新染麹尘罗。\par
章华台畔隋堤上，傍得春风尔许多。\par

\end{document}